% !TeX root = ../../Book3.tex
% 纲要标题：## 第9章　Krull 维数与高度定理
% 纲要位置：工作记录/计划纲要.md，第 2200 行。
% 纲要细目（写作范围）：

\chapter{Krull 维数与高度定理}
\label{b3:ch:09}
\BookChapterNumber{b3:ch:09}

\begin{chapterabstract}
本章用素理想链定义 Krull 维数，比较商环、局部化与支集中的链，并证明主理想定理和高度定理。Noether 正规化把有限型代数的维数与超越次数联系起来；局部维数和极大理想的最少生成元数也得到区分。
\end{chapterabstract}

\begin{learninggoals}
\begin{itemize}
\item 按严格包含的次数计算链长，区分高度、商环维数和模的支集维数。
\item 在主理想定理中辨认极小性，并用高度界限制方程和生成元的个数。
\item 利用超越次数、整扩张及正规化计算有限型代数的维数。
\item 构造参数系，计算嵌入维数，区分参数、极大理想生成元和正则性。
\end{itemize}
\end{learninggoals}

设 \(k\) 为域。在 \(k[x,y]\) 中，可沿 \((0)\subsetneq(x)\subsetneq(x,y)\) 走两步；加上关系 \(xy=0\) 后，\((0)\) 不再是素理想，相应的两个不可约分量各只容许一步素理想链。给环添一个方程，维数往往下降，但下降多少不能只看方程个数：若方程已经由原关系推出，就不会再减少任何链。维数应当从真正存在的素理想链，而非从展示中的符号数目读取。

\ProseParagraphBreak

局部化又只保留落在某个素理想以下的链，因此同一环在不同点附近可能显出不同的局部维数。另一方面，一个局部环的极大理想需要几个生成元，问的是局部坐标有多少独立的一阶方向，与最长素理想链不是同一个问题。比较这两个数字，将为后面对正则点与奇点的区分做好准备。

\ProseParagraphBreak

本章使用\chapref{b3:ch:01}至\chapref{b3:ch:06}中的素谱、局部化、Noether 条件、准素理想、Nakayama 引理和整扩张；超越次数仍指一组超越基的元素个数，其域论基础见第一卷第14.3节。高度定理的证明先于 Hilbert 函数和 Artin--Rees 引理，后面的章节可以直接调用它。正则局部环在这里仅建立定义和基本算例，一般结构定理留在\chapref{b3:ch:18}。

\ProseParagraphBreak
有限型整代数的维数公式可参阅 Matsumura 的《Commutative Ring Theory》\cite[Theorem 5.6, pp. 34--35]{Matsumura1989CommutativeRingTheory}；高度定理和逐次商环的维数见\cite[Theorems 13.5--13.6, pp. 100--101]{Matsumura1989CommutativeRingTheory}，参数系与嵌入维数的定义见\cite[\S14, Systems of parameters and multiplicity, pp. 104--105]{Matsumura1989CommutativeRingTheory}。本章从素理想链证明高度定理。

% !TeX root = ../../Book3.tex
% 纲要标题：### 9.1 维数与高度
% 纲要位置：工作记录/计划纲要.md，第 2079 行。
% 纲要细目（写作范围）：
% * 素理想链、环维数和理想高度
% * 模的维数与支集
% * 局部化对维数的影响

\section{维数与高度}
\label{b3:sec:0901}

素谱中的点既可以独立出现，也可以按包含关系排列。一个不可约闭集内还有多少层更小的不可约闭集，是维数要记录的问题。以下的链长只计算严格包含的次数；不假设所有极大链一样长，也不预设任何几何图像中的维数公式。

\subsection{素理想链、环维数与理想高度}
\label{b3:subsec:0901:01}

\begin{definition}\label{b3:def:krull-dimension-height}
设 \(R\) 为交换环。其有限素理想链
\[
\mathfrak p_0\subsetneq\mathfrak p_1\subsetneq\cdots\subsetneq\mathfrak p_r
\]
的长度 \(r\) 的上确界称为 \(R\) 的\term[Krull-weishu]{Krull 维数}[Krull dimension]，记为 \(\dim R\)。对素理想 \(\mathfrak p\subseteq R\)，以 \(\mathfrak p\) 为末项的此类链长的上确界称为它的\term[gaodu]{高度}[height]，记为 \(\operatorname{ht}\mathfrak p\)。真理想 \(I\subseteq R\) 的高度定义为
\[
\operatorname{ht}I=\inf_{\mathfrak p\supseteq I}\operatorname{ht}\mathfrak p.
\]
允许维数或高度为 \(+\infty\)。零环的谱为空，约定其维数为 \(-\infty\)；单位理想的高度不在本章使用。
\end{definition}
\symidx{\(\dim R\)：环的 Krull 维数}
\symidx{\(\operatorname{ht}\mathfrak p\)、\(\operatorname{ht}I\)：素理想与真理想的高度}

域只有一个素理想 \(0\)，所以维数为零。整数环的非零素理想都是极大理想，故
\(\dim\mathbb Z=1\)。一元多项式环 \(k[t]\) 同样为一维：它是主理想整环，每个非零素理想由不可约多项式生成，因而极大。零维并不表示素谱只有一个点；有限个域的直积已经给出多个点。

\begin{proposition}\label{b3:prop:dimension-basic-comparison}
设 \(R\) 为交换环，\(I\subseteq R\) 为理想，\(R_{\mathrm{red}}=R/\sqrt{(0)}\)。则有 \(\dim(R/I)\le\dim R\)，而
\(\dim R_{\mathrm{red}}=\dim R\)。若 \(\mathfrak p\) 为素理想，且有关数均有限，则
\[
\operatorname{ht}\mathfrak p+\dim(R/\mathfrak p)\le\dim R.
\]
若 \(R_1,\ldots,R_s\) 为交换环且至少一个非零，则
\(\dim(R_1\times\cdots\times R_s)\) 等于非零因子 \(R_i\) 的维数的最大值。
\end{proposition}
\begin{proof}
商环的素理想链对应原环中包含 \(I\) 的链；模掉幂零根时，所有素理想都保留，故链长完全不变。对第二式，取一条末项为 \(\mathfrak p\) 的最长链，以及一条首项为 \(\mathfrak p\) 的最长链，接在一起即可。有限整数的上确界能取到，故这些链存在。直积环的素理想分别来自一个因子，不同因子所对应的素理想互不包含，所以任何链都留在同一个因子中。
\end{proof}

上述不等式不能在没有附加条件时写成等式。高度测量点以下的链，商环维数测量点以上的链；两段能连接起来，不表示任何最长链都经过这个指定点。

\subsection{模的维数与支集}
\label{b3:subsec:0901:02}

\begin{definition}\label{b3:def:module-dimension}
设 \(R\) 为交换环，\(M\) 为 \(R\) 模。支集
\(\operatorname{Supp}_R M=\{\,\mathfrak p\in\operatorname{Spec}R\mid M_{\mathfrak p}\ne0\,\}\)
内有限严格素理想链长的上确界称为 \(M\) 的\term[mo-weishu]{维数}[dimension of a module]，记为 \(\dim_R M\)。空支集的维数取 \(-\infty\)。
\end{definition}
\symidx{\(\dim_R M\)：模的支集维数}

若 \(M\) 有限生成，\cref{b3:thm:support-finite}给出
\(\operatorname{Supp}_R M=V(\operatorname{Ann}_R M)\)，所以
\begin{equation}\label{b3:eq:module-dimension-annihilator}
\dim_R M=\dim\bigl(R/\operatorname{Ann}_R M\bigr).
\end{equation}
这里的维数不是生成元个数，也不是向量空间维数。例如 \(R=k[x]\) 上的 \(R^r\)、\(r\ge1\)，都有支集 \(\operatorname{Spec}R\)，故维数均为一；而 \(R/(x^n)\)、\(n\ge1\)，的支集只有一个闭点，维数为零，其 \(k\) 向量空间维数却为 \(n\)。

\begin{proposition}\label{b3:prop:dimension-associated-components}
若 \(R\) 为交换 Noether 环，\(M\ne0\) 为有限 \(R\) 模，则
\[
\dim_R M=\max_{\mathfrak p\in\operatorname{Ass}_R M}\dim(R/\mathfrak p).
\]
若 \(0\rightarrow L\rightarrow M\rightarrow N\rightarrow0\) 为有限生成 \(R\) 模的短正合列，则
\[
\dim_R M=\max\{\dim_R L,\dim_R N\}.
\]
\end{proposition}
\begin{proof}
由\cref{b3:prop:minimal-primes-associated}，支集中每个素理想都包含某个关联素理想，且关联素理想有限。因此支集等于这些 \(V(\mathfrak p)\) 的有限并。任一链的最小项包含某个这样的 \(\mathfrak p\)，整条链便落在 \(V(\mathfrak p)\) 中，得到第一式。
短正合列的支集是两端支集之并，见\cref{b3:prop:support-exact}。因有限模的支集是闭集，链的最小项落在哪一端，整条链就留在该端，故得到最大值公式。
\end{proof}

\subsection{局部化对维数的影响}
\label{b3:subsec:0901:03}

\begin{proposition}\label{b3:prop:dimension-localization}
设 \(R\) 为交换环，\(\mathfrak p\subseteq R\) 为素理想，\(S\subseteq R\) 为乘法集。则有
\[
\dim R_{\mathfrak p}=\operatorname{ht}\mathfrak p,\qquad
\dim S^{-1}R\le\dim R,\qquad
\dim R=\sup_{\mathfrak m\text{ 极大}}\dim R_{\mathfrak m}.
\]
若 \(M\) 为有限 \(R\) 模，则 \(\dim_{S^{-1}R}S^{-1}M\le\dim_R M\)。零环情形的空上确界同样取 \(-\infty\)。
\end{proposition}
\begin{proof}
局部化素理想对应保留包含关系。\(R_{\mathfrak p}\) 的素理想恰好对应包含于
\(\mathfrak p\) 的素理想，每条这样的链都可在末端补上 \(\mathfrak p\)，所以其维数就是高度。一般局部化只保留不遇到 \(S\) 的素理想，故不能增加链长。任一有限素链的末项包含于某个极大理想；在该极大理想处局部化保留整条链，得到第三式。

对模，若 \(\mathfrak q\cap S=\varnothing\)，则先在 \(S\) 处局部化、再在
\(S^{-1}\mathfrak q\) 处局部化，与直接取 \(M_{\mathfrak q}\) 相同。两者均由允许
\(R\setminus\mathfrak q\) 中的分母给出，分式的互逆公式保证这一识别。因此局部模的支集链对应原支集中避开 \(S\) 的链，得到不等式。
\end{proof}

一般点处的局部化可能大幅降低维数。例如整环的零理想处局部环是分式域，维数为零，不论原环有多少维。相反，在多项式环的闭点处，局部化会保留许多不同方向的素理想链；下一节将用理想生成元的个数限制这些链的长度。

% !TeX root = ../../Book3.tex
% 纲要标题：### 9.2 高度定理
% 纲要位置：工作记录/计划纲要.md，第 2085 行。
% 纲要细目（写作范围）：
% * Krull 主理想定理
% * 广义高度定理
% * 生成元数与余维约束

\section{高度定理}
\label{b3:sec:0902}

一个方程不应一次切掉任意多层素理想链，这个直觉在 Noether 环中由主理想定理保证。证明需要小心区分极小素理想与任意包含该方程的素理想。本节先用准素性和 Nakayama 引理建立单个方程的结论，再归纳处理多个方程。

\subsection{Krull 主理想定理}
\label{b3:subsec:0902:01}

\begin{lemma}\label{b3:lem:principal-primary-local-domain}
设 \((R,\mathfrak m)\) 为 Noether 局部整环，\(a\in R\)，且
\(\sqrt{aR}=\mathfrak m\)。则 \(\dim R\le1\)。
\end{lemma}
\begin{proof}
若 \(a=0\)，整环性给出 \(\mathfrak m=0\)，故 \(R\) 是域。以下设 \(a\ne0\)。
反设存在非零素理想 \(0\ne\mathfrak p\subsetneq\mathfrak m\)。因
\(\sqrt{aR}=\mathfrak m\)，有 \(a\notin\mathfrak p\)。

环 \(R/aR\) 是零维 Noether 环，所以由\cref{b3:thm:noether-dimension-zero}为 Artin 环。考虑符号幂
\(Q_n=\mathfrak p^nR_{\mathfrak p}\cap R\)。在 \(R_{\mathfrak p}\) 中，\(\mathfrak p^nR_{\mathfrak p}\) 的根为唯一极大理想，故由\cref{b3:prop:maximal-radical-primary}准素；再由\cref{b3:prop:primary-localization}，收缩后的每个 \(Q_n\) 都是 \(\mathfrak p\) 准素理想。它们在 \(R/aR\) 中的像组成降链，因此存在 \(n\ge1\) 使
\[
Q_n+aR=Q_{n+1}+aR.
\]
对 \(x\in Q_n\)，可写 \(x=y+ar\)，其中 \(y\in Q_{n+1}\)。于是 \(ar\in Q_n\)；
因 \(a\notin\sqrt{Q_n}=\mathfrak p\)，准素性迫使 \(r\in Q_n\)。所以
\[
Q_n=Q_{n+1}+aQ_n.
\]
有限生成模 \(Q_n/Q_{n+1}\) 因而等于自身的 \(a\) 倍。由 \(a\in\mathfrak m\) 及
\cref{b3:thm:nakayama}，它为零，故 \(Q_n=Q_{n+1}\)。

在 \(\mathfrak p\) 处局部化得到
\(\mathfrak p^nR_{\mathfrak p}=\mathfrak p^{n+1}R_{\mathfrak p}\)。
再对 Noether 局部环 \(R_{\mathfrak p}\) 上的有限生成模
\(\mathfrak p^nR_{\mathfrak p}\) 使用 Nakayama 引理，得到它为零。但
\(R_{\mathfrak p}\) 是整环，且非零的 \(\mathfrak p\) 在其中仍非零，其正整数次幂不可能为零，矛盾。于是除零理想外，没有严格小于 \(\mathfrak m\) 的素理想。
\end{proof}

\begin{theorem}[Krull 主理想定理]\label{b3:thm:principal-ideal-theorem}
设 \(R\) 为交换 Noether 环，\(a\in R\)。每个在 \((a)\) 上极小的素理想 \(\mathfrak p\) 都满足
\[
\operatorname{ht}\mathfrak p\le1.
\]
\end{theorem}
\begin{proof}
在 \(\mathfrak p\) 处局部化后，只存在一个包含 \(a\) 的素理想，即唯一极大理想；
因此可设 \(R\) 为局部环且 \(\sqrt{aR}=\mathfrak m\)。若存在长度至少二的素链，可截取并在末端补上 \(\mathfrak m\)，得到
\(\mathfrak q_0\subsetneq\mathfrak q_1\subsetneq\mathfrak m\)。
在局部整环 \(R/\mathfrak q_0\) 中，\(a\) 的像的根仍为极大理想的像，上一引理却排除了非零的中间素理想
\(\mathfrak q_1/\mathfrak q_0\)，矛盾。因此局部维数至多为一，即所求高度界。
\end{proof}

这称为\term[Krull-zhulixiang-dingli]{Krull 主理想定理}[Krull's principal ideal theorem]。
在整环中，若 \(a\) 非零且非单位，其上极小素理想非零，故高度恰好为一。
“极小”不可删除：例如 \((x)\subseteq(x,y)\) 于 \(k[x,y]\)，后者包含主理想
\((x)\)，却有链 \(0\subsetneq(x)\subsetneq(x,y)\)，高度至少为二。

\subsection{广义高度定理}
\label{b3:subsec:0902:02}

\begin{theorem}[广义高度定理]\label{b3:thm:height-theorem}
若 \(R\) 为交换 Noether 环，\(I=(a_1,\ldots,a_r)\subsetneq R\)，其中 \(r\ge0\)，则每个在 \(I\) 上极小的素理想 \(\mathfrak p\) 都满足
\[
\operatorname{ht}\mathfrak p\le r.
\]
\end{theorem}
\begin{proof}
对 \(r\) 归纳。\(r=0\) 时，\(\mathfrak p\) 是极小素理想，高度为零；\(r=1\) 已由主理想定理证明。一般情形在 \(\mathfrak p\) 处局部化，可设
\((R,\mathfrak m)\) 局部且 \(\sqrt I=\mathfrak m\)。

取任意链
\[
\mathfrak p_0\subsetneq\cdots\subsetneq\mathfrak p_d=\mathfrak m.
\]
若 \(d=0\) 无须证明。因 \(I\not\subseteq\mathfrak p_{d-1}\)，重排生成元后可设
\(a_1\notin\mathfrak p_{d-1}\)。令 \(\mathfrak q_d=\mathfrak m\)，从
\(i=d-1\) 向下逐次在 \(\mathfrak q_{i+1}\) 内选一个在
\(\mathfrak p_i+(a_1)\) 上极小的素理想 \(\mathfrak q_i\)。这样的选择可以先在
\(\mathfrak q_{i+1}\) 处局部化、取一个极小素理想，再收缩；任何更小的候选也在
\(\mathfrak q_{i+1}\) 内，故它在原环中仍然极小。

对 \(i\le d-2\)，若 \(\mathfrak q_i=\mathfrak q_{i+1}\)，则
\[
\mathfrak p_i\subsetneq\mathfrak p_{i+1}\subsetneq\mathfrak q_i.
\]
第二个包含严格，因为 \(a_1\in\mathfrak q_i\) 而
\(a_1\notin\mathfrak p_{i+1}\)。但在 \(R/\mathfrak p_i\) 中，
\(\mathfrak q_i/\mathfrak p_i\) 在一个主理想上极小，主理想定理禁止它具有这样的长度二链。故
\(\mathfrak q_0\subsetneq\cdots\subsetneq\mathfrak q_{d-1}\) 严格。

这是一条包含 \(a_1\) 的长度 \(d-1\) 的链。在局部商环 \(R/(a_1)\) 中，
其极大理想是剩余 \(r-1\) 个生成元的根，由归纳假设高度至多 \(r-1\)。
所以 \(d-1\le r-1\)，即 \(d\le r\)。原链任意，结论成立。
\end{proof}

这一定理也称\term[Krull-gaodu-dingli]{Krull 高度定理}[Krull's height theorem]。
证明没有将
\(\operatorname{ht}\mathfrak p\) 擅自拆成两个高度之和；一般环上的素链未必具有允许这种拆分的长度性质。Matsumura 的《Commutative Ring Theory》给出相同高度结论\cite[Theorem 13.5, pp. 100--101]{Matsumura1989CommutativeRingTheory}，但其证明次序先用 Samuel 函数；本卷采用上述独立论证，使后面的增长理论能够调用高度定理。

\subsection{生成元数与余维约束}
\label{b3:subsec:0902:03}

若 \((R,\mathfrak m)\) 为 Noether 局部环，\(\mathfrak m\) 有有限个生成元，故
\cref{b3:thm:height-theorem}说明 \(\dim R<\infty\)。更一般地，若
\(\sqrt{(a_1,\ldots,a_r)}=\mathfrak m\)，就有 \(\dim R\le r\)。
这限制的是使零点集中到闭点所需的方程数；不表示任一主理想的任一零点都只有余维一。

\begin{lemma}\label{b3:lem:dimension-cutting}
设 \((R,\mathfrak m)\ne0\) 为交换 Noether 局部环，\(d=\dim R\)。存在
\(x_1,\ldots,x_d\in\mathfrak m\)，使
\[
\sqrt{(x_1,\ldots,x_d)}=\mathfrak m.
\]
\(d=0\) 时取空组。
\end{lemma}
\begin{proof}
对有限整数 \(d\) 归纳。\(d=0\) 时唯一素理想为 \(\mathfrak m\)，所以幂零根为
\(\mathfrak m\)，空组满足要求。若 \(d>0\)，环的极小素理想只有有限多个，且都严格小于
\(\mathfrak m\)；否则 \(\mathfrak m\) 极小便排除了其他素理想，维数将为零。
由\cref{b3:lem:finite-prime-avoidance}选
\(x_1\in\mathfrak m\) 避开全部极小素理想。

任一包含 \(x_1\) 的素链，其首项不是极小素理想，可以在下面补上一项。
因此 \(\dim R/(x_1)\le d-1\)。对这个局部商环使用归纳假设，提升其至多
\(d-1\) 个元素，得到一个根为 \(\mathfrak m\)、至多由 \(d\) 个元素生成的理想。
若少于 \(d\) 个，则广义高度定理会给出 \(\dim R<d\)，矛盾。所以总数恰为 \(d\)，完成归纳。
\end{proof}

这个引理与高度定理合起来，将局部维数解释为生成一个 \(\mathfrak m\) 准素理想所需的最少元素个数。\secref{b3:sec:0904}将把这样的元素组命名为参数系。它们生成的理想通常小于 \(\mathfrak m\)，因此参数个数和极大理想的最少生成元数仍需区分。

% !TeX root = ../../Book3.tex
% 纲要标题：### 9.3 有限型代数的维数
% 纲要位置：工作记录/计划纲要.md，第 2214 行。
% 纲要细目（写作范围）：
% * 多项式环的维数
% * Noether 正规化和超越次数
% * 整扩张的维数保持

\section{有限型代数的维数}
\label{b3:sec:0903}

多项式环中有明显的素理想链，但仅展示一条长链只能给出下界。上界需要证明：在有限型整代数中，每次模掉一个非零素理想，分式域的超越次数就会下降。Noether 正规化随后把一般有限型整代数与多项式环联系起来，借整性将多项式环中的长链提升到原环。


\subsection{多项式环的维数}
\label{b3:subsec:0903:01}

取 \(A=k[x,y]\)、\(\mathfrak p=(x)\)，商环 \(A/\mathfrak p=k[\bar y]\) 的分式域的一个超越基是 \(\{\bar y\}\)。把 \(\bar y\) 提升为 \(y\)，再添入被取商杀掉的非零元素 \(x\)，就在 \(\operatorname{Frac}(A)\) 中得到代数独立的 \(x,y\)。下面的引理说明，在一般有限型整代数中，从商环的有限代数生成组中选出的超越基提升后，也能这样添入一个非零素理想中的元素；困难在于证明添入后仍然代数独立。

\begin{lemma}\label{b3:lem:transcendence-degree-drop}
设 \(A\) 为域 \(k\) 上的有限型整代数，\(0\ne\mathfrak p\) 为素理想。则
\[
\operatorname{trdeg}_k\operatorname{Frac}(A/\mathfrak p)
<\operatorname{trdeg}_k\operatorname{Frac}(A).
\]
\end{lemma}
\begin{proof}
从 \(A/\mathfrak p\) 的有限代数生成组中取极大代数独立组
\(\bar y_1,\ldots,\bar y_r\)，并在 \(A\) 中取提升 \(y_i\)。其他生成元在
\(k(\bar y_1,\ldots,\bar y_r)\) 上代数，故左侧超越次数是 \(r\)。提升后的 \(y_i\) 也代数独立。

取 \(0\ne a\in\mathfrak p\)。我们证明 \(a\) 在 \(k(y_1,\ldots,y_r)\) 上超越。
若它代数，取它在域 \(K=k(y_1,\ldots,y_r)\) 上的首一最小多项式
\(\mu(T)\in K[T]\)。为了在模掉 \(\mathfrak p\) 后留下关于 \(\bar y_i\) 的非零关系，需要这条关系有一个不含 \(a\) 的非零常数项。
若 \(\mu(0)=0\)，则 \(T\) 整除不可约首一多项式 \(\mu\)，故只能有 \(\mu(T)=T\)，从而 \(a=0\)，与所选元素非零矛盾。因此 \(\mu(0)\ne0\)。
选取 \(\mu\) 各系数的非零公共分母 \(D(y)\in k[y_1,\ldots,y_r]\)，
将 \(\mu(a)=0\) 乘以 \(D(y)\)，得到
\[
c_s(y)\cdot a^s+\cdots+c_1(y)\cdot a+c_0(y)=0,
\qquad 0\ne c_0\in k[Y_1,\ldots,Y_r].
\]
这里 \(D(y)\mu(0)\ne0\)，因为 \(D(y)\) 和 \(\mu(0)\) 均为域 \(K\) 中的非零元；
\(y_1,\ldots,y_r\) 的代数独立性又保证对应的常数项多项式 \(c_0(Y)\) 非零。
模掉 \(\mathfrak p\) 得
\(c_0(\bar y)=0\)，与代数独立性矛盾。因此 \(y_1,\ldots,y_r,a\) 代数独立，右侧超越次数至少是 \(r+1\)。
\end{proof}

\begin{theorem}\label{b3:thm:polynomial-dimension}
任意域 \(k\) 及整数 \(n\ge0\) 满足
\[
\dim k[X_1,\ldots,X_n]=n.
\]
更一般地，有限型整 \(k\) 代数 \(A\) 的维数不超过
\(\operatorname{trdeg}_k\operatorname{Frac}(A)\)。
\end{theorem}
\begin{proof}
对严格素链 \(\mathfrak p_0\subsetneq\cdots\subsetneq\mathfrak p_r\)，逐次将上一引理用于
\(A/\mathfrak p_i\) 及其非零素理想 \(\mathfrak p_{i+1}/\mathfrak p_i\)。严格包含使这个新加入的素关系非零，因此超越次数每一步至少下降一。一条长度 \(r\) 的素链就需要起点至少有 \(r\) 个代数独立参数；而起点的超越次数不超过 \(\operatorname{trdeg}_k\operatorname{Frac}(A)\)，因为商环中的代数独立元素提升后仍然代数独立。这证明上界。对于多项式环，
\[
(0)\subsetneq(X_1)\subsetneq(X_1,X_2)\subsetneq\cdots
\subsetneq(X_1,\ldots,X_n)
\]
的各项都是素理想，给出长度 \(n\) 的链。\(n=0\) 时环为域，结论同样成立。
\end{proof}

例如 \(k[X,Y,Z]/(XY)\) 的两个极小素理想为 \((\bar X)\) 和 \((\bar Y)\)，相应商环均为二元多项式环，故其维数为二。这个计算按不可约分量分别进行，不把有零因子的环擅自当成具有一个分式域的整环。

\subsection{Noether 正规化与超越次数}
\label{b3:subsec:0903:02}

\begin{theorem}\label{b3:thm:affine-dimension}\label{b3:thm:affine-dimension-trdeg}
设 \(k\) 为域。若 \(A\) 是有限型整 \(k\) 代数，则
\[
\dim A=\operatorname{trdeg}_k\operatorname{Frac}(A).
\]
若 \(A\ne0\) 只是有限型交换 \(k\) 代数，则
\[
\dim A=\max_{\mathfrak p\in\operatorname{Min}A}
\operatorname{trdeg}_k\operatorname{Frac}(A/\mathfrak p).
\]
这里 \(\operatorname{Min}A\) 表示极小素理想的集合。
\end{theorem}
\begin{proof}
整环情形由\cref{b3:thm:noether-normalization}选出代数独立元素
\(y_1,\ldots,y_d\)，使 \(A\) 是 \(B=k[y_1,\ldots,y_d]\) 上的有限生成模。
\secref{b3:sec:0504}已经证明，这种有限整包含给出有限域扩张
\(k(y_1,\ldots,y_d)\subseteq\operatorname{Frac}(A)\)。因此这些 \(y_i\) 构成该分式域的超越基，超越次数为 \(d\)。由\cref{b3:thm:polynomial-dimension}，这给出上界 \(\dim A\le d\)。

下界来自 \(B\subseteq A\) 的整性。多项式环 \(B\) 中有长度 \(d\) 的素链，由\cref{b3:thm:lying-over}提升首项，再反复用\cref{b3:thm:going-up}提升其余各项；收缩严格增加，提升链也严格增加。因此 \(\dim A\ge d\)，两界合起来得到维数公式。

一般情形下，\(A\) 为 Noether 环，只有有限多个极小素理想。任意素链的首项都包含某个极小素理想，因此整条链都来自某个 \(A/\mathfrak p\)。反之这些商环中的链也是 \(A\) 中的链。这给出
\(\dim A=\max_{\mathfrak p\in\operatorname{Min}A}\dim(A/\mathfrak p)\)，再应用整环情形。
\end{proof}

一般情形的公式说明，全环维数由维数最大的不可约分量决定，不同分量可以有不同维数。幂零厚度不改变素理想链；嵌入素理想所对应的闭集已包含在某个不可约分量中，也不会增加这个最大值。

尖点环 \(k[t^2,t^3]\) 整包含 \(k[t^2]\)，其分式域为 \(k(t)\)，所以维数为一。
变量之间有关系 \((t^3)^2=(t^2)^3\)，两个代数生成元不意味着二维。
相反，\(k[x,y,x^{-1}]\) 虽添入了一个逆元，分式域仍为 \(k(x,y)\)，所以维数仍为二。
有限型条件在这里保证 Noether 正规化可用；不能把该公式未经说明用于任意子环。

\subsection{整扩张的维数保持}
\label{b3:subsec:0903:03}

\begin{theorem}\label{b3:thm:integral-dimension}
若交换环间的包含 \(A\subseteq B\) 是整扩张，则 \(\dim A=\dim B\)。不要求扩张有限，也不要求两环为整环。
\end{theorem}
\begin{proof}
取 \(B\) 中的严格素链 \(\mathfrak q_0\subsetneq\cdots\subsetneq\mathfrak q_r\)，令 \(\mathfrak p_i=\mathfrak q_i\cap A\)。这些收缩仍然严格：若相邻两项相等，\cref{b3:thm:incomparability}就会迫使相应的 \(\mathfrak q_i\) 相等。因此收缩保留链长，得到 \(\dim B\le\dim A\)。

反之，从 \(A\) 中任意有限严格素链出发，先用\cref{b3:thm:lying-over}提升首项，再用\cref{b3:thm:going-up}逐项提升。每个提升的收缩是原来指定的素理想，不同的收缩保证每次包含严格。所以提升也保留链长，得到 \(\dim B\ge\dim A\)。这比较的是所有有限链的长度，维数无穷时也成立。
\end{proof}

整包含 \(A\subseteq B\) 使两个素谱的整体维数相同，但这不能直接给出一个指定素理想 \(\mathfrak q\) 与其收缩的高度相等。Going-up 提升链时，上端的素理想由构造产生，未必等于预先指定的 \(\mathfrak q\)。若要固定这个上端，逐项向下补出原像，就需要 Going-down；当基环 \(A\) 为整闭整环、扩环 \(B\) 也为整环时，它给出下面的点处结论。
\begin{corollary}\label{b3:cor:integral-height-normal-base}
若 \(A\subseteq B\) 为整环间的整扩张，且 \(A\) 整闭，则对
\(\mathfrak q\in\operatorname{Spec}B\)、\(\mathfrak p=\mathfrak q\cap A\)，有
\(\operatorname{ht}\mathfrak q=\operatorname{ht}\mathfrak p\)。
\end{corollary}
\begin{proof}
以 \(\mathfrak q\) 为末项的链收缩后仍严格，故 \(\operatorname{ht}\mathfrak q\le\operatorname{ht}\mathfrak p\)。反向由\cref{b3:thm:going-down}，从指定的 \(\mathfrak q\) 开始向下提升每条以 \(\mathfrak p\) 为末项的有限链。不同的收缩使提升链严格，于是 \(\operatorname{ht}\mathfrak p\le\operatorname{ht}\mathfrak q\)。
\end{proof}

整同态若不是单射，应先把基环换成其像：此时能推出的是
\[
\dim B=\dim\bigl(A/\ker(A\rightarrow B)\bigr).
\]
例如商同态
\(k[x]\rightarrow k\)、\(x\mapsto0\) 是整同态，却并不保持原基环的维数。

% !TeX root = ../../Book3.tex
% 纲要标题：### 9.4 局部维数、参数系与嵌入维数
% 纲要位置：工作记录/计划纲要.md，第 2097 行。
% 纲要细目（写作范围）：
% * 参数系
% * 嵌入维数与 Krull 维数
% * 正则局部环的定义先行，结构理论留给第18章
% ---


\section{局部维数、参数系与嵌入维数}
\label{b3:sec:0904}

在 Noether 局部环中，维数有两种容易混淆的比较对象：使零点只剩闭点需要多少个方程，以及生成整个极大理想需要多少个元素。前者等于 Krull 维数，后者可能更大。两者相等时得到正则局部环的定义。


\subsection{参数系}
\label{b3:subsec:0904:01}

\begin{definition}\label{b3:def:system-of-parameters}
设 \((R,\mathfrak m)\ne0\) 为交换 Noether 局部环，\(d=\dim R\)。若一组元素
\(x_1,\ldots,x_d\in\mathfrak m\) 满足
\(\sqrt{(x_1,\ldots,x_d)}=\mathfrak m\)，则称它们为 \(R\) 的\term[canshuxi]{参数系}[system of parameters]，称它们生成的理想为\term[canshu-lixiang]{参数理想}[parameter ideal]。
\end{definition}

\begin{theorem}\label{b3:thm:system-of-parameters}
设 \((R,\mathfrak m)\ne0\) 为交换 Noether 局部环，\(d=\dim R\)。则 \(R\) 有参数系，且 \(d\) 等于生成一个根为 \(\mathfrak m\) 的理想所需的最少元素个数。
若 \(x_1,\ldots,x_d\) 是参数系，则对 \(0\le i\le d\)，有
\[
\dim R/(x_1,\ldots,x_i)=d-i.
\]
此外，对任意 \(a_1,\ldots,a_r\in\mathfrak m\)，都有
\(\dim R/(a_1,\ldots,a_r)\ge d-r\)。
\end{theorem}
\begin{proof}
存在性已由\cref{b3:lem:dimension-cutting}证明，高度定理说明少于 \(d\) 个元素不可能生成根为 \(\mathfrak m\) 的理想。
设 \(e=\dim R/(a_1,\ldots,a_r)\)，在此商环中选一个有 \(e\) 个元素的参数系并提升。
它与 \(a_1,\ldots,a_r\) 合起来生成根为 \(\mathfrak m\) 的理想，所以 \(d\le r+e\)。
对参数系的前 \(i\) 项应用这个不等式，得到商维数至少是 \(d-i\)；剩余 \(d-i\) 项在商环中生成极大根理想，高度定理又给出反向不等式。
\end{proof}

\(R/(x_1,\ldots,x_d)\) 是零维 Noether 局部环，故是 Artin 环，且作为 \(R\) 模具有有限长度。
反之，若商模具有有限长度，其支集只含闭点，于是生成理想的根为 \(\mathfrak m\)。所以参数系也可以通过商环的有限长度识别。

\ProseParagraphBreak
例如在 \(k[x,y]_{(x,y)}\) 中，\(x,y\) 和 \(x^2,y^3\) 都是参数系，后者并不生成极大理想。
在尖点局部环 \(k[t^2,t^3]_{(t^2,t^3)}\) 中，\(t^2\) 单独构成参数系：模掉它后
\((t^3)^2=(t^2)^3=0\)，故商环只有一个素理想。参数并不要求能作自由的坐标。

\subsection{嵌入维数与 Krull 维数}
\label{b3:subsec:0904:02}

\begin{definition}\label{b3:def:embedding-dimension}
设 \((R,\mathfrak m)\ne0\) 为交换 Noether 局部环，剩余域为 \(\kappa=R/\mathfrak m\)。向量空间 \(\mathfrak m/\mathfrak m^2\) 在 \(\kappa\) 上的维数
\[
\dim_\kappa(\mathfrak m/\mathfrak m^2)
\]
称为 \(R\) 的\term[qianru-weishu]{嵌入维数}[embedding dimension]，记为 \(\operatorname{edim}R\)。
\end{definition}
\symidx{\(\operatorname{edim}R\)：局部环的嵌入维数}

\(\mathfrak m/\mathfrak m^2\) 上的标量作用来自 \(R\)，并因 \(\mathfrak m\) 零化该商模而降为 \(\kappa\) 作用。
\cref{b3:cor:local-generators}说明，嵌入维数恰好等于极大理想的最少生成元数。
由高度定理立即得到
\begin{equation}\label{b3:eq:dimension-embedding-bound}
\dim R\le\operatorname{edim}R.
\end{equation}

这个向量空间仅保留函数的一次信息；乘积落入 \(\mathfrak m^2\) 而消失。取
\[
R=k[x_1,\ldots,x_n]_{(x_1,\ldots,x_n)},
\]
则变量的像构成
\(\mathfrak m/\mathfrak m^2\) 的基。生成性显然；若线性式属于局部化后的平方理想，则乘一个常数项非零的分母后，其一次部分仍是该线性式乘以非零常数，故各系数必须为零。因此嵌入维数为 \(n\)。坐标素链在此局部化中全部保留，Krull 维数也为 \(n\)。

\ProseParagraphBreak
尖点的情形不同。\(t^2,t^3\) 在 \(\mathfrak m/\mathfrak m^2\) 中线性独立，因为平方理想中每个非零元素的最低次数至少为四，局部化的单位分母不改变这一结论。因此该环的嵌入维数为二，Krull 维数为一。这两个数字之差已经能辨认这个点与直线上的点不同。

\subsection{正则局部环的定义}
\label{b3:subsec:0904:03}

\begin{definition}\label{b3:def:regular-local-ring}
若非零交换 Noether 局部环 \((R,\mathfrak m)\) 满足
\(\dim R=\dim_{R/\mathfrak m}(\mathfrak m/\mathfrak m^2)\)，则称它为
\term[zhengze-jubuhuan]{正则局部环}[regular local ring]。
\end{definition}

域是零维正则局部环；反过来，零维正则局部环的极大理想满足
\(\mathfrak m/\mathfrak m^2=0\)，由 Nakayama 引理有 \(\mathfrak m=0\)，所以它是域。
上小节的多项式局部环是正则的，尖点局部环则不是。
另一个简单的非正则例子为 \(k[\epsilon]/(\epsilon^2)\)：其 Krull 维数为零，而嵌入维数为一。

\ProseParagraphBreak
设 \((R,\mathfrak m)\) 为 \(d\) 维正则局部环。极大理想 \(\mathfrak m\) 的任一极小生成组
都有 \(d\) 项，并且生成的理想就是 \(\mathfrak m\)，因而构成参数系；这样的生成组称为\term[zhengze-canshuxi]{正则参数系}[regular system of parameters]。
正则序列在\chapref{b3:ch:15}引入，正则参数系构成正则序列的证明见\secref{b3:sec:1801}。
正规表示整闭，正则表示上述维数等式，它们是不同定义。下一章将在一维整环情形研究二者与离散赋值环的关系。

\begin{table}[htbp]
\centering
\caption[维数类不变量]{维数类不变量。\(R\) 为交换含幺环，\(M\) 为 \(R\) 模；局部环行中的 \(\mathfrak m\) 表示唯一极大理想，\(\kappa=R/\mathfrak m\)。空支集与零环的维数取 \(-\infty\)。}
\label{b3:tab:dimension-invariants}
\begin{tabularx}{\textwidth}{@{}>{\raggedright\arraybackslash}p{0.19\textwidth}>{\raggedright\arraybackslash}p{0.23\textwidth}>{\raggedright\arraybackslash}X@{}}
\toprule
不变量 & 条件 & 所测量的对象与结论\\
\midrule
\(\operatorname{ht}\mathfrak p\) & \(\mathfrak p\in\operatorname{Spec}R\) & 以 \(\mathfrak p\) 为末项的严格素链长的上确界，等于 \(\dim R_{\mathfrak p}\)\\
\(\dim(R/\mathfrak p)\) & \(\mathfrak p\in\operatorname{Spec}R\) & 以 \(\mathfrak p\) 为首项的严格素链长的上确界，即闭子空间 \(V(\mathfrak p)\) 的维数\\
\(\dim_R M\) & 任意 \(R\) 模 \(M\) & 支集 \(\operatorname{Supp}_R M\) 内的严格素链长的上确界；若 \(M\) 有限生成，等于 \(\dim(R/\operatorname{Ann}_R M)\)\\
\(\operatorname{edim}R\) & \((R,\mathfrak m)\ne0\) 为 Noether 局部环 & \(\dim_\kappa(\mathfrak m/\mathfrak m^2)\)，等于生成 \(\mathfrak m\) 所需的最少元素个数\\
生成 \(\mathfrak m\) 准素理想的最少元素数 & \((R,\mathfrak m)\ne0\) 为 Noether 局部环 & 在所有 \(\mathfrak m\) 准素理想的生成组中取元素个数的最小值，等于 \(\dim R\)；达到最小值的生成组是参数系\\
\bottomrule
\end{tabularx}
\end{table}

\cref{b3:tab:dimension-invariants}中的高度测量指定素理想以下的链，商环维数测量它以上的链。若各量有限，\cref{b3:prop:dimension-basic-comparison}给出
\[
\operatorname{ht}\mathfrak p+\dim(R/\mathfrak p)\le\dim R.
\]
\begin{samepage}
取一条经过 \(\mathfrak p\)、以极大理想 \(\mathfrak m\) 为末项的素链，可以在 \(\mathfrak p\) 处分成两段：
\[
\underbrace{\mathfrak p_0\subsetneq\cdots\subsetneq\mathfrak p}
_{u\le\operatorname{ht}\mathfrak p}
\underbrace{\subsetneq\cdots\subsetneq\mathfrak m}
_{v\le\dim(R/\mathfrak p)}.
\]
这里 \(u,v\) 是两段各自的长度。对两段分别取最长链再拼接，就得到上述不等式；原环的最长链未必经过这个指定素理想。
\end{samepage}


\begin{chaptersummary}
素链给出维数的定义，也给出最直接的下界。上界需要额外论证：Noether 条件下由理想的生成元数控制高度，有限型整代数中则由超越次数沿严格素链的下降控制链长。Lying-over 提供素链首项的提升，Going-up 逐项向上提升，Incomparability 保证收缩后的包含仍然严格，从而整扩张保持整体维数；指定点的高度比较另外用到 Going-down 的假设。

\ProseParagraphBreak
在 Noether 局部环中，参数系存在且长度恰为维数，但它生成的理想通常不是极大理想。极大理想最少需要多少个生成元，由剩余域向量空间 \(\mathfrak m/\mathfrak m^2\) 决定。尖点的 Krull 维数为一，而嵌入维数为二，说明这两种计数回答的是不同问题。
\end{chaptersummary}

\BookExercises

\begin{exercise}\label{b3:ex:09-product-dimension}
设 \(k\) 为域。对任意整数 \(d\ge2\) 和 \(0\le h\le d-2\)，构造一个 \(k\) 上的 Noether 代数 \(R\) 及其素理想 \(\mathfrak p\)，使
\[
\dim R=d,\qquad \operatorname{ht}\mathfrak p=h,\qquad
\dim(R/\mathfrak p)=1.
\]
由此说明高度与商环维数之和可以比全环维数小任意给定的正整数。
\end{exercise}
\begin{solution}
取
\[
C=k[x_1,\ldots,x_h,t],\qquad B=k[y_1,\ldots,y_d],\qquad
R=C\times B,\qquad \mathfrak p=(x_1,\ldots,x_h)C\times B.
\]
两个多项式环都 Noether，故其有限直积也 Noether。由\cref{b3:thm:polynomial-dimension}和直积维数公式，
\(\dim R=\max\{h+1,d\}=d\)。商环 \(R/\mathfrak p\cong k[t]\)，所以 \(\mathfrak p\) 素且商维数为一。
直积中的素链只留在同一个因子内，故 \(\operatorname{ht}_R\mathfrak p\) 等于 \((x_1,\ldots,x_h)\) 在 \(C\) 中的高度。逐个加入 \(x_i\) 得到长度 \(h\) 的素链，而高度定理给出反向上界，因而该高度恰为 \(h\)；\(h=0\) 时这个理想是整环 \(C\) 的零理想，高度为零。
差值为 \(d-h-1\)。给定正整数 \(g\)，取 \(h=0,d=g+1\)，便得到差值 \(g\)。
\end{solution}

\begin{exercise}\label{b3:ex:09-reduced-dimension}
设 \(k\) 为域，\(A=k[x,y,\epsilon]/(\epsilon^3)\)。计算 \(\dim A\) 和 \(\dim_A A/(x,\epsilon)\)，并解释增加幂零元为何没有增加维数。
\end{exercise}
\begin{solution}
每个素理想都含 \(\epsilon\)，故 \(A\) 的谱与 \(k[x,y]\) 有相同素链，维数二。
商模是 \(k[y]\)，其支集对应商环的谱，模维数一。幂零元给出不同的环乘法信息，却不产生新的素理想或链。
\end{solution}

\begin{exercise}\label{b3:ex:09-nonfinite-support}
将 \(\mathbb Q/\mathbb Z\) 视为 \(\mathbb Z\) 模，计算其零化子、支集和模维数，说明有限生成模的零化子公式不能任意推广。
\end{exercise}
\begin{solution}
没有非零整数零化全部有理数类，故零化子为零。在零素理想处局部化，每个元素原本被某个非零整数零化，所以局部化为零。
在 \((p)\) 处，\(1/p+\mathbb Z\) 不能被不属于 \((p)\) 的整数零化，因此局部化非零。
支集恰为全部非零素理想，这些素理想互不可比，模维数为零；而
\(\dim\bigl(\mathbb Z/\operatorname{Ann}(\mathbb Q/\mathbb Z)\bigr)=1\)。
\end{solution}

\begin{exercise}\label{b3:ex:09-support-sum}
设 \(k\) 为域，\(R=k[x,y,z]\)。计算 \(R/(x)\oplus R/(y,z)\) 及子模
\(xR\subseteq R\) 的维数。再由短正合列核对 \(\dim R\)。
\end{exercise}
\begin{solution}
两直和项的维数为二和一，直和维数为二。乘以 \(x\) 给出 \(R\simeq xR\)，故后者维数三。
短正合列 \(0\rightarrow xR\rightarrow R\rightarrow R/(x)\rightarrow0\) 的维数公式为
\(3=\max\{3,2\}\)。非零子模不必比商模低维。
\end{solution}

\begin{exercise}\label{b3:ex:09-local-coordinate-height}
设 \(k\) 为域，\(n\ge0\) 为整数，\(R=k[x_1,\ldots,x_n]\)。令 \(\mathfrak p=(x_1,\ldots,x_r)\)、\(0\le r\le n\)。
计算 \(\dim R_{\mathfrak p}\) 和 \(\dim R/\mathfrak p\)。
\end{exercise}
\begin{solution}
前 \(r\) 个变量逐个加入给出末项为 \(\mathfrak p\) 的长度 \(r\) 素链。
另一方面 \(\mathfrak p\) 在自己生成的 \(r\) 元理想上极小，高度定理给出高度至多 \(r\)。所以
\(\dim R_{\mathfrak p}=\operatorname{ht}\mathfrak p=r\)，而商环为 \(n-r\) 元多项式环，维数为 \(n-r\)。
\end{solution}

\begin{exercise}\label{b3:ex:09-height-components}
设 \(k\) 为域，\(R=k[x,y,z]\)。求 \(I=(xy,xz)\) 的极小素理想及其高度，并计算 \(\dim R/I\)。
\end{exercise}
\begin{solution}
包含 \(I\) 的素理想若不含 \(x\)，必同时含 \(y,z\)，所以极小素理想恰为
\((x)\) 与 \((y,z)\)，高度分别为一和二。这符合两个生成元的高度上界，却不要求所有极小素理想有相同高度。
两个不可约分量的坐标环维数分别为二和一，因此 \(\dim R/I=2\)。
\end{solution}

\begin{exercise}\label{b3:ex:09-primary-generators}
设 \(k\) 为域，\(R=k[x,y]_{(x,y)}\)，\(a,b\) 为正整数。证明任何 \((x,y)R\) 准素理想都不能由一个元素生成。
再求 \(R/(x^a,y^b)\) 的长度。
\end{exercise}
\begin{solution}
该局部环的维数为二。若一个主理想的根为极大理想，主理想定理会给出维数至多一，矛盾。
所给商环中 \(x^iy^j\)、\(0\le i<a,0\le j<b\)，构成 \(k\) 基；常数项非零的分母在这个环中已经可逆，因为其非常数部分幂零。商环为局部 Artin 环且剩余域为 \(k\)，每个简单因子是一维 \(k\) 空间，所以长度为 \(ab\)。
\end{solution}

\begin{exercise}\label{b3:ex:09-hypersurface-dimension}
设 \(k\) 为域，\(n\ge1\) 为整数，\(f\in k[x_1,\ldots,x_n]\) 为非零非单位。证明
\(\dim k[x_1,\ldots,x_n]/(f)=n-1\)。
\end{exercise}
\begin{solution}
先设 \(f\) 不可约，重排变量使其对 \(x_n\) 的次数为正。它不能整除只含前
\(n-1\) 个变量的非零多项式，所以这些变量在商环中仍代数独立。
关系 \(f=0\) 又说明 \(x_n\) 在它们的分式域上代数；首项系数虽未必为单位，在分式域中却可逆。
因此商整环的分式域超越次数为 \(n-1\)，用有限型维数公式得结论。
一般 \(f\) 的极小素理想由其不可约因子生成；有限多个对应商环都有维数 \(n-1\)，故结论不变。
\end{solution}

\begin{exercise}\label{b3:ex:09-open-dimension}
设 \(k\) 为域，\(A\) 为有限型整 \(k\) 代数，\(0\ne f\in A\)。证明 \(\dim A_f=\dim A\)。
解释这为何不与 \(\dim A_{(0)}=0\) 冲突。
\end{exercise}
\begin{solution}
只倒置 \(f\) 可写为 \(A_f\cong A[u]/(fu-1)\)，所以 \(A_f\) 仍是有限型整 \(k\) 代数。它与 \(A\) 具有相同分式域，由\cref{b3:thm:affine-dimension}得到两者维数相等。
零素理想处的局部化则倒置全部非零元素，得到 \(A_{(0)}=\operatorname{Frac}(A)\)。若 \(\dim A>0\)，这个域不可能是有限型 \(k\) 代数：否则由 Zariski 引理\cref{b3:lem:zariski}，它在 \(k\) 上有限代数，从而超越次数为零，与 \(\dim A>0\) 矛盾。因而不能对它继续使用有限型整代数的维数公式。例如 \(k[t]_{(0)}=k(t)\) 的环维数为零，超越次数却为一。
\end{solution}

\begin{exercise}\label{b3:ex:09-monomial-curve-dimension}
设 \(k\) 为域，\(t\) 为不定元，\(a,b\) 为互素正整数。计算 \(k[t^a,t^b]\) 的维数和分式域，并说明哪一步使用互素性。
\end{exercise}
\begin{solution}
整数 Bézout 等式 \(ra+sb=1\) 给出 \(t=(t^a)^r(t^b)^s\)，允许负指数后说明分式域为 \(k(t)\)，所以维数一。
互素性只用于认出分式域等于 \(k(t)\)。若最大公因数为 \(d\)，相同论证给出分式域 \(k(t^d)\)，超越次数仍为一，因此维数仍为一。
\end{solution}

\begin{exercise}\label{b3:ex:09-nonfinite-integral}
设 \(L/k\) 为无限次代数扩张。证明 \(k[t]\subseteq L[t]\) 为非有限的整扩张，并计算两环维数。
\end{exercise}
\begin{solution}
任意一个 \(L[t]\) 元素只有有限多个系数，它们生成有限域扩张 \(k'/k\)。环
\(k'[t]\) 是 \(k[t]\) 上有限自由模，所以其中每个元素都整。这证明整性。
若 \(L[t]\) 是有限生成的 \(k[t]\) 模，模掉 \(t\) 后会使 \(L\) 成为有限维 \(k\) 空间，矛盾。
两环分别是一元多项式环，维数均为一，符合整扩张的维数不变性。
\end{solution}

\begin{exercise}\label{b3:ex:09-local-integral-height}
设 \(k\) 为域，\(A=k[x,y]\)，\(B=k[x,y,z]/(z^2-xy)\)。证明 \(A\) 自然嵌入 \(B\)，且 \(B\) 是整环。计算 \(B\) 在素理想 \((x,z)\) 和 \((x,y,z)\) 处的局部维数，并说明高度比较为什么需要把素链的末项固定。
\end{exercise}
\begin{solution}
关于 \(z\) 的首一关系允许带余除法，每个商类都唯一写成 \(a+bz\)，其中 \(a,b\in A\)。因此 \(B\cong A\oplus Az\) 作为 \(A\) 模，常数商类使 \(A\rightarrow B\) 单射；\(B\) 又是有限生成 \(A\) 模，故该包含为有限整扩张。

整环性另由关系的不可约性得到。多项式 \(z^2-xy\) 对 \(k[x,y]\) 中的素元 \(x\) 满足 Eisenstein 判据\footnote{艾森斯坦（Ferdinand Gotthold Max Eisenstein，1823-1852），德国数学家，研究数论与椭圆函数。}：常数项被 \(x\) 整除但不被 \(x^2\) 整除，其余非首项系数也被 \(x\) 整除。因此它不可约，唯一分解环 \(k[x,y,z]\) 中的不可约元为素元，故商环 \(B\) 是整环。
两理想的商环分别为 \(k[y]\) 和 \(k\)，所以确为素理想；收缩到 \(A\) 后分别是 \((x)\) 和 \((x,y)\)。基环 \(A\) 为整闭整环，\cref{b3:cor:integral-height-normal-base}使两处高度分别等于一和二，这也就是所求局部维数。
Going-up 从已选的底端素理想向上构造链，末项未必是这两个指定素理想之一；Going-down 则能从指定末项向下提升。这解释了所用高度比较中整闭基环的作用。
\end{solution}

\begin{exercise}\label{b3:ex:09-parameter-zero-divisor}
设 \(k\) 为域，\(n\ge2\) 为整数，\(R=\bigl(k[x,y]/(x^2,xy^n)\bigr)_{(x,y)}\)。计算 \(R\) 的 Krull 维数、嵌入维数及 \(\operatorname{Ann}_R(y)\)。再对 \(a,b\in k\)，判定 \(ax+by\) 何时构成参数，并证明满足条件的这些参数全是零因子。
\end{exercise}
\begin{solution}
商环中的每个元素唯一写成 \(f(y)+xg(y)\)，其中 \(f\in k[y]\)，\(g\in k[y]/(y^n)\)。局部化后相应地有
\[
R=k[y]_{(y)}\oplus x\bigl(k[y]/(y^n)\bigr)
\]
作为 \(k[y]_{(y)}\) 模的分解：常数项非零的分母在第二项上也可逆，而第二项的平方为零。幂零根因此为 \((x)\)，约化局部环为 \(k[y]_{(y)}\)，故 \(\dim R=1\)。原关系都没有一次项，\(x,y\) 的像在 \(\mathfrak m/\mathfrak m^2\) 中构成一组基，故 \(\operatorname{edim}R=2\)。

若 \(y(f+xg)=0\)，上面的直和分解使 \(yf=0\) 及 \(yg=0\) 在各自的项中成立。第一项是整环，故 \(f=0\)；第二项中 \(yg=0\) 当且仅当 \(g\in k y^{n-1}\)。所以
\[
\operatorname{Ann}_R(y)=(xy^{n-1}),\qquad xy^{n-1}\ne0.
\]
元素 \(ax+by\) 在约化环中的像是 \(by\)。若 \(b=0\)，该元素幂零，生成理想的根为 \((x)\)，不等于 \(\mathfrak m\)。若 \(b\ne0\)，商环中可令 \(y=-ab^{-1}x\)，得到 \(R/(ax+by)\cong k[x]/(x^2)\)，其维数为零，所以该元素是参数。此时
\((ax+by)xy^{n-1}=ax^2y^{n-1}+bxy^n=0\)，被零化的 \(xy^{n-1}\) 非零，故这个参数是零因子。
\end{solution}

\begin{exercise}\label{b3:ex:09-parameter-powers}
设 \((R,\mathfrak m)\ne0\) 为交换 Noether 局部环，\(d=\dim R\)，\(x_1,\ldots,x_d\) 为 \(R\) 的参数系。证明其中每个元素可独立换为任意正整数次幂，仍得参数系。又证明它们可以任意重排，且每个前 \(i\) 项的商环维数为 \(d-i\)（\(0\le i\le d\)）。
\end{exercise}
\begin{solution}
一个素理想包含元素的正整数次幂当且仅当包含元素本身，因此这些理想有相同根。
重排不改变生成理想。对重排后的参数系应用\cref{b3:thm:system-of-parameters}，便得每个前缀的维数公式。
\end{solution}

\begin{exercise}\label{b3:ex:09-artin-regular}
设 \(R\ne0\) 为交换 Artin 局部环。证明 \(R\) 正则当且仅当它是域。再设 \(k\) 为域、\(n\ge0\) 为整数，计算
\(k[x_1,\ldots,x_n]/(x_1,\ldots,x_n)^2\) 的嵌入维数。
\end{exercise}
\begin{solution}
Artin 环 Noether 且零维。正则性等价于 \(\mathfrak m/\mathfrak m^2=0\)，Nakayama 引理使 \(\mathfrak m=0\)，故为域；域显然正则。
所给环的极大理想平方为零，\(x_1,\ldots,x_n\) 的像构成其 \(k\) 基，所以嵌入维数为 \(n\)。当 \(n>0\) 时它不正则。
\end{solution}

\begin{exercise}\label{b3:ex:09-tangent-linear-relations}
设 \(k\) 为域，\(n\ge1\) 为整数。令 \(\mathfrak n=(x_1,\ldots,x_n)\subseteq k[x_1,\ldots,x_n]\)，并设
\(I=(f_1,\ldots,f_r)\subseteq\mathfrak n\)，\(R=\bigl(k[x_1,\ldots,x_n]/I\bigr)_{\mathfrak n/I}\)。
证明 \(\operatorname{edim}R\) 等于 \(n\) 减去这些 \(f_i\) 的一次部分的秩。
\end{exercise}
\begin{solution}
极大理想的平方商为 \(\mathfrak n/(\mathfrak n^2+I)\)，局部化不改变这个被 \(\mathfrak n\) 零化的 \(k\) 向量空间。
模掉 \(\mathfrak n^2\) 后，每个 \(f_i\) 只留下其一次部分；任意组合
\(\sum_i g_if_i\) 的一次部分为 \(\sum_i g_i(0)(f_i)_{\mathrm{lin}}\)。
所以全部关系恰是这些线性形式的张成空间，商空间维数为题述数值。
\end{solution}

\begin{exercise}\label{b3:ex:09-nonpure-parameter}
设 \(k\) 为域，\(R=\bigl(k[x,y,z]/(xy,xz)\bigr)_{(x,y,z)}\)。证明其维数为二、嵌入维数为三，并验证 \(x+y,z\) 是参数系。
\end{exercise}
\begin{solution}
两个极小素理想对应的商环为局部化后的 \(k[y,z]\) 和 \(k[x]\)，维数分别为二和一，所以总维数二。
所定义关系都为二次式，没有一次关系，嵌入维数三。
令 \(z=0,y=-x\) 后，商环为 \(k[x]_{(x)}/(x^2)\)，是零维局部环，故两个元素构成参数系。
这个例子也说明参数系的存在不要求各个不可约分量等维。
\end{solution}

\begin{exercise}\label{b3:ex:09-module-parameters}
设 \((R,\mathfrak m)\) 为交换 Noether 局部环，\(M\ne0\) 为有限生成的 \(R\) 模，\(d=\dim_R M\)。
证明存在 \(x_1,\ldots,x_d\in\mathfrak m\) 使 \(M/(x_1,\ldots,x_d)M\) 具有有限长度，并证明用 \(\mathfrak m\) 中少于 \(d\) 个元素不能做到这一点。
\end{exercise}
\begin{solution}
由\cref{b3:thm:support-finite}，\(\operatorname{Supp}_R M=V(\operatorname{Ann}_R M)\)，所以 \(A=R/\operatorname{Ann}_R M\) 的素谱恰好记录 \(M\) 的支集，且 \(\dim A=d\)。对 \(A\) 取参数系，相当于在这个支集上选出足够多的方程。将参数提升到 \(R\)，得到理想
\(I=(x_1,\ldots,x_d)\)。有限生成模的支集及局部 Nakayama 引理说明
\[
\operatorname{Supp}(M/IM)=\operatorname{Supp}(M)\cap V(I).
\]
具体地，局部化后若 \(I\not\subseteq\mathfrak p\)，商为零；若 \(I\subseteq\mathfrak p\) 且
\(M_{\mathfrak p}\ne0\)，Nakayama 引理禁止其等于 \(I M_{\mathfrak p}\)，所以商非零。
令 \(Q=M/IM\)。参数条件使 \(\operatorname{Supp}_R Q=\{\mathfrak m\}\)；\(Q\) 有限生成，故由支集的零化子公式得
\[
\sqrt{\operatorname{Ann}_R Q}=\mathfrak m.
\]
这意味着 \(S=R/\operatorname{Ann}_R Q\) 是只有一个素理想的 Noether 环，故由\cref{b3:thm:noether-dimension-zero}为 Artin 环。现在 \(Q\) 是有限生成的 \(S\) 模，\cref{b3:thm:artin-finite-length}给出其有限长度。这里调用的证明正是利用幂零极大理想的有限滤过：\(\mathfrak m^NQ=0\) 对某个 \(N\) 成立，且各商层 \(\mathfrak m^iQ/\mathfrak m^{i+1}Q\) 是有限维剩余域向量空间。\(R\) 与 \(S\) 在 \(Q\) 上有相同的子模，因此这也是它作为 \(R\) 模的长度。

反之，若 \(J=(y_1,\ldots,y_r)\subseteq\mathfrak m\) 使 \(M/JM\) 有有限长度，Nakayama 引理保证这个商非零。其支集因此为 \(\{\mathfrak m\}\)。由上述支集等式，\(J\) 在 \(A\) 中的像的根是 \(\mathfrak m A\)；高度定理给出 \(d=\dim A\le r\)，所以少于 \(d\) 个元素不能做到这一点。
\end{solution}
